A few facts about photometry of asteroids:
\begin{itemize}
\item Asteroids are small, irregularly shaped objects of our Solar System
  that orbit the Sun in approximately elliptical orbits.
\item Their brightness observed at a given instant from Earth depend on the
  surface area illuminated by the Sun and the part of the asteroid which is
  visible to the observer. Both vary as the asteroid moves.
\item The way the sunlight is reflected by the surface of the asteroid
  depends on its texture and on the angle between the Sun, the asteroid and
  the observer (phase angle), which varies as the Earth and the asteroid move
  along their orbits. In particular, asteroids with surfaces covered by fine
  dust (called regoliths) exhibit a sharp increase in brightness at phase
  angles $\varphi$ close to zero (i.e., when they are close to the
  opposition).
\item Since the observed flux of any source decreases with the square of
  distance, the observed magnitude of an asteroid also depends on its
  distance from the Sun and from the observer at the time of the observation.
  Their apparent magnitude $m$ outside the atmosphere is then
  \[ M(t) = M_r(t) + 5\log(RD) \]
  where $m_r$ is usually called reduced magnitude (meaning the magnitude the
  asteroid would have if its distances from the Sun and the Earth were
  reduced both to 1\,AU) and depends only on the visible illuminated surface
  area and on phase angle effects. $R$ and $D$ are, respectivelly, the
  heliocentric and geocentric distances.
\end{itemize}

Consider now the following scenario. Light curves of a given asteroid were
obtained at three different nights at different points of its orbit, and at
each time a photometric standard was observed in the same frame of the
asteroid. Table 1 shows the geometric configuration of the asteroid at each
night (phase angle $\varphi$, in degrees, $R$ and $D$ in AUs), and the
calibrated magnitude of the photometric standard star that was observed along
with the asteroid. Consider the calibrated magnitude as the final apparent
magnitude after all the effects are taken into account.

Tables 2, 3 and 4 contain, for each night, the time interval of each
observation with respect to the first one (in hours) the air mass, the
instrumental magnitude of the asteroid and the instrumental magnitude of the
star.

Air mass is the dimensionless thickness of the atmosphere along the line of
sight, and is equal to 1 when looking to zenith.

\begin{center}
Table 1

\begin{tabular}{c|c|c|c|c}
Night & $D$ & $R$ & $\varphi$ & $M_{\mathrm{star}}$ \\
\hline
A & 0.36 & 1.35 & 0.0  & 8.2 \\
B & 1.15 & 2.13 & 8.6  & 8.0 \\
C & 2.70 & 1.89 & 15.6 & 8.1 \\
\end{tabular}
\end{center}

\begin{center}
Table 2: Night A

\begin{tabular}{c|c|c|c}
$\Delta t$ & Air mass & $m_{\mathrm{ast}}$ & $m_{\mathrm{star}}$ \\
\hline
0.00 & 1.28 & 7.44 & 8.67 \\
0.44 & 1.18 & 7.38 & 8.62 \\
0.89 & 1.11 & 7.34 & 8.59 \\
1.33 & 1.06 & 7.28 & 8.58 \\
1.77 & 1.02 & 7.32 & 8.58 \\
2.21 & 1.00 & 7.33 & 8.56 \\
2.66 & 1.00 & 7.33 & 8.56 \\
3.10 & 1.01 & 7.30 & 8.56 \\
3.54 & 1.03 & 7.27 & 8.58 \\
3.99 & 1.07 & 7.27 & 8.58 \\
4.43 & 1.13 & 7.31 & 8.61 \\
4.87 & 1.21 & 7.37 & 8.63 \\
5.31 & 1.32 & 7.42 & 8.67 \\
5.76 & 1.48 & 7.49 & 8.73 \\
6.20 & 1.71 & 7.59 & 8.81 \\
6.64 & 2.06 & 7.69 & 8.92 \\
7.09 & 2.62 & 7.87 & 9.14 \\
7.53 & 3.67 & 8.21 & 9.49 \\
\end{tabular}
\end{center}

\begin{center}
Table 3: Night B

\begin{tabular}{c|c|c|c}
$\Delta t$ & Air mass & $m_{\mathrm{ast}}$ & $m_{\mathrm{star}}$ \\
\hline
0.00 & 1.28 & 13.24 & 8.38 \\
0.44 & 1.18 & 13.21 & 8.36 \\
0.89 & 1.11 & 13.13 & 8.34 \\
1.33 & 1.06 & 13.11 & 8.33 \\
1.77 & 1.02 & 13.11 & 8.32 \\
2.21 & 1.00 & 13.15 & 8.32 \\
2.66 & 1.00 & 13.17 & 8.32 \\
3.10 & 1.01 & 13.17 & 8.32 \\
3.54 & 1.03 & 13.13 & 8.33 \\
3.99 & 1.07 & 13.15 & 8.34 \\
4.43 & 1.13 & 13.14 & 8.34 \\
4.87 & 1.21 & 13.14 & 8.37 \\
5.31 & 1.32 & 13.21 & 8.38 \\
5.76 & 1.48 & 13.30 & 8.43 \\
6.20 & 1.71 & 13.34 & 8.47 \\
6.64 & 2.06 & 13.39 & 8.54 \\
7.09 & 2.62 & 13.44 & 8.65 \\
7.53 & 3.67 & 13.67 & 8.87 \\
\end{tabular}
\end{center}

\begin{center}
Table 4: Night C

\begin{tabular}{c|c|c|c}
$\Delta t$ & Air mass & $m_{\mathrm{ast}}$ & $m_{\mathrm{star}}$ \\
\hline
0.00 & 1.28 & 11.64 & 8.58 \\
0.44 & 1.18 & 11.53 & 8.54 \\
0.89 & 1.11 & 11.56 & 8.60 \\
1.33 & 1.06 & 11.49 & 8.52 \\
1.77 & 1.02 & 11.58 & 8.48 \\
2.21 & 1.00 & 11.79 & 8.63 \\
2.66 & 1.00 & 11.67 & 8.53 \\
3.10 & 1.01 & 11.53 & 8.46 \\
3.54 & 1.03 & 11.47 & 8.48 \\
3.99 & 1.07 & 11.63 & 8.67 \\
4.43 & 1.13 & 11.51 & 8.51 \\
4.87 & 1.21 & 11.65 & 8.55 \\
5.31 & 1.32 & 11.77 & 8.61 \\
5.76 & 1.48 & 11.88 & 8.75 \\
6.20 & 1.71 & 11.86 & 8.78 \\
6.64 & 2.06 & 12.03 & 9.03 \\
7.09 & 2.62 & 12.14 & 9.19 \\
7.53 & 3.67 & 12.63 & 9.65 \\
\end{tabular}
\end{center}

\begin{enumerate}
\item Plot the star magnitude versus air mass for each set
\item Calculate the extinction coefficient for each night (see Appendix A at
  the end of the text)
\item Were the observations affected by clouds in one night? Express your
  answer as: a) Night A, b) Night B, c) Night C, d) None of the nights.
\item Plot the calibrated magnitude versus time for each set of observations
  of the asteroid (see Appendix B).
\item Determine the rotation period for each night. Consider that the light
  curve for this asteroid has two minima and two maxima, and that the
  semi-period is the average of the intervals between the two maxima and the
  two minima.
\item Determine the amplitude (difference from maximum to minimum) of the
  light curve for each night
\item Plot the calibrated reduced magnitude $M_r$ versus phase angle
  $\varphi$ (use the mean value of each light curve)
\item Calculate the angular coefficient of the phase curve (the plot of the
  calibrated reduced magnitude versus the phase angle) considering only the
  points away from the opposition (See 3rd bullet of facts about photometry
  of asteroids above)
\item Is there any reason to assume a surface covered by fine dust (regolith)?
  Answer YES/NO.
\end{enumerate}

\textbf{Appendix A: extinction coefficients}

The magnitude outside the atmosphere is given by:
\[ M = m - A.X + B \]
where $A$ is the extinction coefficient and $B$ is the zero point coefficient
for the night, $X$ is the airmass of the observation and $m$ is the
instrumental magnitude (that is, the magnitude obtained directly from the
image)

\textbf{Appendix B: calibrated differential magnitude}

If the standard star is the same frame as the object, the calibrated
magnitude of the object can be obtained by:
\[ M_{ast} = m_{ast} - m_{star} + M_{star} \]

\textbf{Appendix C: least squares estimation of angular coefficients}

Given straight lines described by
\[ y_i = \alpha + \beta x_i, \quad i = 1..n \]
the least squares estimator of $\beta$ is given by
\[ \beta = \frac{\sum_i^n x_i y_i
     - \frac{1}{n} \sum_i^n x_i \sum_i^n y_i}
   {\sum_i^n x_i^2 - \frac{1}{n} \left( \sum_i^n x_i \right)^2} \]
